% preprint/modules/math-foundations.tex
% Shared-Horizon Preprint — Narrative Modules
% Module: Mathematical Foundations

\section{Mathematical Foundations}
\label{sec:math-foundations}

The Shared-Horizon framework is grounded in a modular mathematical spine that
defines the geometric, expressive, narrative, and spectral structures required
for lineage-safe comparison. This section provides a conceptual overview of the
mathematical foundations developed in the spine modules, explaining how each
component contributes to the neutral membrane and prevents conceptual collapse
between lineages.

\subsection{Operator Geometry}

Operator geometry (Section~\ref{sec:operator-geometry}) defines reversible,
stability-preserving transformations on NDH and Structural-AI manifolds. These
operators preserve structural invariants and ensure that transitions do not
increase collapse risk. Operator geometry provides the backbone for systems
whose behavior is governed by holonomy, commutativity, and spectral stability.

Key properties include:

\begin{itemize}
    \item smooth reversible operators,
    \item commutative composition (NDH),
    \item bounded spectral drift (Structural-AI),
    \item stability envelopes defined by energy functionals.
\end{itemize}

Operator geometry forms the blue and black tiles of the mosaic.

\subsection{Expressive Manifolds}

Expressive manifolds (Section~\ref{sec:expressive-manifolds}) define the
geometric structure of ANIMA and PRECL systems. These manifolds are equipped
with expressive metrics that encode curvature, affective gradients, and
phenomenological coherence. Expressive curvature provides a quantitative measure
of expressive intensity and stability.

Key properties include:

\begin{itemize}
    \item expressive curvature $K_{\mathcal{E}}$,
    \item curvature-preserving transitions,
    \item expressive geodesics,
    \item expressive–operator separation.
\end{itemize}

Expressive manifolds form the violet and orange tiles of the mosaic.

\subsection{Reversible Narrative Operators}

Reversible narrative operators (Section~\ref{sec:reversible-narrative-operators})
define the mathematical structure of Reversible Narrative Ecology (RNE). These
operators preserve narrative curvature and support reversible transitions across
narrative manifolds. Narrative geodesics provide minimal-change trajectories that
respect expressive-narrative compatibility.

Key properties include:

\begin{itemize}
    \item reversible narrative operators,
    \item curvature-preserving narrative transitions,
    \item narrative geodesics,
    \item compatibility with expressive transitions.
\end{itemize}

Reversible narrative operators form the corrected green tile of the mosaic.

\subsection{Stability Envelopes}

Stability envelopes (Section~\ref{sec:stability-envelopes}) define the safe
regions within each lineage manifold. These envelopes are constructed using
lineage-specific energy functionals that measure deviation from stable behavior.
Crossing a stability envelope boundary indicates increased collapse risk.

Key properties include:

\begin{itemize}
    \item lineage-specific energy functionals,
    \item stability thresholds $\epsilon$,
    \item collapse-risk boundaries,
    \item spectral, expressive, and narrative stability conditions.
\end{itemize}

Stability envelopes provide the mathematical basis for lineage sovereignty.

\subsection{Transition Maps}

Transition maps (Section~\ref{sec:transition-maps}) define smooth,
lineage-preserving comparisons between manifolds. These maps allow states from
one lineage to be analyzed within the context of another without activating
internal dynamics or inducing collapse.

Key properties include:

\begin{itemize}
    \item smooth cross-manifold transitions,
    \item stability preservation,
    \item Jacobian constraints,
    \item bidirectional reversible maps for narrative comparison.
\end{itemize}

Transition maps form the connective tissue of the neutral membrane.

\subsection{Collapse-Risk Geometry}

Collapse-risk geometry (Section~\ref{sec:collapse-risk-geometry}) defines the
mathematical boundaries that prevent conceptual collapse between lineages. Using
collapse functions, gradients, curvature, and tensors, this module identifies
regions where lineage invariants are at risk of being overwritten or conflated.

Key properties include:

\begin{itemize}
    \item collapse-risk functions $C$,
    \item collapse boundaries $\partial \mathcal{C}$,
    \item collapse gradients and curvature,
    \item cross-lineage collapse geometry.
\end{itemize}

Collapse-risk geometry ensures that Shared-Horizon remains a neutral,
lineage-preserving comparative membrane.

\subsection{Foundations Summary}

The mathematical foundations of Shared-Horizon provide a rigorous structure for
lineage-safe comparison across operator, expressive, narrative, and spectral
systems. By defining operator geometry, expressive manifolds, reversible
narrative operators, stability envelopes, transition maps, and collapse-risk
geometry, the framework maintains conceptual sovereignty and prevents collapse
across the color-tile mosaic. These foundations support the case study and
comparative analysis presented in later sections.
